\exnumber{06}
\extitle{UX VI - Ex 7 - Itération}
\exfiles{ex06/archive zip, nomme Name\_FirstName\_DeliverableName\_Date\_VersionNumber}
\exauthorize{Ceux nommés dans l'exercice}
\makeheaderfiles

\textbf{Instructions :}

À partir de vos tests précédents, choisissez l'un de vos 3 prototypes.

Implémentez les 3 idées que vous avez notées comme suggestions.

Réalisez une autre série de tests (et adaptez votre protocole de test si nécessaire).

Rédigez un rapport avec les conclusions que vous avez tirées de cette boucle d'itération.

\textbf{Livrables :}

\begin{itemize}
\item Photos du nouveau test
\item Un rapport avec vos conclusions, vos impressions et vos apprentissages clés de cette itération (document PDF issu de Microsoft Word, Pages, Google Docs ou tout logiciel de traitement de texte)
\item Comparaison entre la V1 et la V2 du prototype dans le rapport
\item Déposez les documents sur Github, dans un zip nommé « ProjectTitle\_Name\_FirstName », et chaque livrable ainsi : « Name\_FirstName\_DeliverableName\_Date\_VersionNumber ». Exemple : « Doe\_John\_MidFiPrototype\_05242024\_V1 ».
\item Ajoutez un titre dans tous les documents écrits.
\item Soignez l'aspect global des livrables. Ils doivent être clairs, simples et lisibles. Vous pouvez regarder canva.com, sans surcharger le dessin de détails.
\end{itemize}

\textbf{Pistes :}

\begin{itemize}
\item Ce rapport sera remis au CEO. Soyez attentif aux détails et assurez-vous que le document est propre et facile à utiliser !
\end{itemize}

